% GENERATED FILE. Edit canonical lessons, then run npm run generate:books.
% canonical-source-hash: fnv1a64:a128c4db198f017e
% canonical-lessons: SA-C114-mesah, SA-C114-kapotah, SA-C114-ulukah, SA-C114-bakah, SA-C114-bhramarah

\chapter{Sheep, Pigeon, Owl, Heron, Bee}
\label{ch:sa-sheep-pigeon-owl-heron-bee}
\hlchaptermodality{114}

\begin{chapteropening}
\textbf{By the end of this chapter:} \emph{I can say a sheep, a pigeon, an owl, a heron and a bee.}

Complete the last lesson of chapter 114: I can say a sheep, a pigeon, an owl, a heron and a bee.
\end{chapteropening}

\section[\texorpdfstring{\sk{मेषः} (meṣaḥ)}{मेषः (meṣaḥ)}]{\sk{मेषः} (meṣaḥ) — a sheep}

\label{lesson:SA-C114-mesah}



\begin{quote}
Before the new one: say the Sanskrit for a bull, then the Sanskrit for a buffalo.
\end{quote}

\subsection*{You'll want to know: \sk{मेषः}}
\textbf{\sk{मेषः}} — \emph{meṣaḥ} — \textquotedblleft{}a sheep\textquotedblright{}.

\begin{grammarlens}[title={What you've built}]
One more word for this chapter.
\end{grammarlens}

\subsection*{Your turn}
\begin{itemize}
\raggedright
  \item \emph{Say it:} \emph{meṣaḥ}
  \item \emph{Say it:} \emph{meṣaḥ}, once more
  \item \emph{Say it:} say \emph{meṣaḥ}
  \item \emph{Recall:} say the Sanskrit for a bull, then the Sanskrit for a buffalo, then say \emph{meṣaḥ} again
\end{itemize}

\begin{tcolorbox}[breakable,colback=teal!4,colframe=teal!35!black,title={Before you move on}]
What does \sk{मेषः} mean? (\textquotedblleft{}A sheep\textquotedblright{}.) Say it once more.
\end{tcolorbox}
\section[\texorpdfstring{\sk{कपोतः} (kapotaḥ)}{कपोतः (kapotaḥ)}]{\sk{कपोतः} (kapotaḥ) — a pigeon}

\label{lesson:SA-C114-kapotah}



\begin{quote}
Before the new one: say the Sanskrit for a buffalo, then the Sanskrit for a sheep.
\end{quote}

\subsection*{You'll want to know: \sk{कपोतः}}
\textbf{\sk{कपोतः}} — \emph{kapotaḥ} — \textquotedblleft{}a pigeon\textquotedblright{}.

\begin{grammarlens}[title={What you've built}]
One more word for this chapter.
\end{grammarlens}

\subsection*{Your turn}
\begin{itemize}
\raggedright
  \item \emph{Say it:} \emph{kapotaḥ}
  \item \emph{Say it:} \emph{kapotaḥ}, once more
  \item \emph{Say it:} say \emph{kapotaḥ}
  \item \emph{Recall:} say the Sanskrit for a buffalo, then the Sanskrit for a sheep, then say \emph{kapotaḥ} again
\end{itemize}

\begin{tcolorbox}[breakable,colback=teal!4,colframe=teal!35!black,title={Before you move on}]
What does \sk{कपोतः} mean? (\textquotedblleft{}A pigeon\textquotedblright{}.) Say it once more.
\end{tcolorbox}
\section[\texorpdfstring{\sk{उलूकः} (ulūkaḥ)}{उलूकः (ulūkaḥ)}]{\sk{उलूकः} (ulūkaḥ) — an owl}

\label{lesson:SA-C114-ulukah}



\begin{quote}
Before the new one: say the Sanskrit for a sheep, then the Sanskrit for a pigeon.
\end{quote}

\subsection*{You'll want to know: \sk{उलूकः}}
\textbf{\sk{उलूकः}} — \emph{ulūkaḥ} — \textquotedblleft{}an owl\textquotedblright{}.

\begin{grammarlens}[title={What you've built}]
One more word for this chapter.
\end{grammarlens}

\subsection*{Your turn}
\begin{itemize}
\raggedright
  \item \emph{Say it:} \emph{ulūkaḥ}
  \item \emph{Say it:} \emph{ulūkaḥ}, once more
  \item \emph{Say it:} say \emph{ulūkaḥ}
  \item \emph{Recall:} say the Sanskrit for a sheep, then the Sanskrit for a pigeon, then say \emph{ulūkaḥ} again
\end{itemize}

\begin{tcolorbox}[breakable,colback=teal!4,colframe=teal!35!black,title={Before you move on}]
What does \sk{उलूकः} mean? (\textquotedblleft{}An owl\textquotedblright{}.) Say it once more.
\end{tcolorbox}
\section[\texorpdfstring{\sk{बकः} (bakaḥ)}{बकः (bakaḥ)}]{\sk{बकः} (bakaḥ) — a heron}

\label{lesson:SA-C114-bakah}



\begin{quote}
Before the new one: say the Sanskrit for a pigeon, then the Sanskrit for an owl.
\end{quote}

\subsection*{You'll want to know: \sk{बकः}}
\textbf{\sk{बकः}} — \emph{bakaḥ} — \textquotedblleft{}a heron\textquotedblright{}.

\begin{grammarlens}[title={What you've built}]
One more word for this chapter.
\end{grammarlens}

\subsection*{Your turn}
\begin{itemize}
\raggedright
  \item \emph{Say it:} \emph{bakaḥ}
  \item \emph{Say it:} \emph{bakaḥ}, once more
  \item \emph{Say it:} say \emph{bakaḥ}
  \item \emph{Recall:} say the Sanskrit for a pigeon, then the Sanskrit for an owl, then say \emph{bakaḥ} again
\end{itemize}

\begin{tcolorbox}[breakable,colback=teal!4,colframe=teal!35!black,title={Before you move on}]
What does \sk{बकः} mean? (\textquotedblleft{}A heron\textquotedblright{}.) Say it once more.
\end{tcolorbox}
\section[\texorpdfstring{\sk{भ्रमरः} (bhramaraḥ)}{भ्रमरः (bhramaraḥ)}]{\sk{भ्रमरः} (bhramaraḥ) — a bee}

\label{lesson:SA-C114-bhramarah}



\begin{quote}
Before the new one: say the Sanskrit for an owl, then the Sanskrit for a heron.
\end{quote}

\subsection*{You'll want to know: \sk{भ्रमरः}}
\textbf{\sk{भ्रमरः}} — \emph{bhramaraḥ} — \textquotedblleft{}a bee\textquotedblright{}.

\begin{grammarlens}[title={What you've built}]
One more word for this chapter.
\end{grammarlens}

\subsection*{Your turn}
\begin{itemize}
\raggedright
  \item \emph{Say it:} \emph{bhramaraḥ}
  \item \emph{Say it:} \emph{bhramaraḥ}, once more
  \item \emph{Say it:} say \emph{bhramaraḥ}
  \item \emph{Recall:} say the Sanskrit for an owl, then the Sanskrit for a heron, then say \emph{bhramaraḥ} again
\end{itemize}

\begin{tcolorbox}[breakable,colback=teal!4,colframe=teal!35!black,title={Before you move on}]
What does \sk{भ्रमरः} mean? (\textquotedblleft{}A bee\textquotedblright{}.) Say it once more.
\end{tcolorbox}
